\documentclass[10pt,a4paper]{amsart}
\usepackage[margin=19mm]{geometry}
\usepackage{amsmath,amssymb,mathtools}
\usepackage[scr=rsfs,cal=euler]{mathalfa}
\usepackage{xfrac}
\usepackage[unicode,hidelinks,bookmarks=false]{hyperref}
\newcommand{\ie}{i.e.\ }
\newcommand{\cf}{cf.\ }
\newcommand{\resp}{respectively\ }
\newcommand{\Subsection}{\subsection}
\providecommand{\nobreakdash}{\nobreak\hbox{-}}
\setlength{\emergencystretch}{2em}
\multlinegap0pt
\usepackage{amssymb}
\usepackage{mathtools}
\usepackage{tikz}
\usepackage[shortlabels]{enumitem}
\usepackage{xcolor}
\newtheorem{lem}{Lemma}
\newtheorem{notation}{Notation}
\newtheorem{prop}{Proposition}
\usepackage{cleveref}
\crefname{lem}{Lemma}{Lemmas}
\crefname{notation}{Notation}{Notations}
\crefname{prop}{Proposition}{Propositions}
\title{Algebraic Invariants of Edge Ideals Under Suspension}
\author{Selvi Kara, Dalena Vien}
\date{Source: arXiv:2603.05657v1 (CC BY 4.0)}
\begin{document}
\maketitle

\noindent\emph{Source and licence.} Algebraic Invariants of Edge Ideals Under Suspension. Version arXiv:2603.05657v1 (CC BY 4.0), distributed by arXiv under the Creative Commons Attribution 4.0 International licence, http://creativecommons.org/licenses/by/4.0/, as recorded in the arXiv licence metadata for this exact version and shown on the abstract page https://arxiv.org/abs/2603.05657v1. Full licence notice: \texttt{licenses/arxiv-2603.05657-CC-BY-4.0.txt}.\par\smallskip

\noindent\textit{Editorial scope.} Counted unit: the article's prop prop:inclusion-nonzero (source0.tex lines 928--939). The article's complete original proof of it is delivered with it (source0.tex lines 941--994, one \texttt{proof} environment of 54 lines). No mathematics is written, repaired or re-derived: the statement and the proof are the article's own lines, in the article's own notation, and no retained line is substituted or re-flowed.  Retained uncounted as the source-local context the proof needs: lines 361--371; lines 464--471; lines 880--898; lines 903--913.  Nothing was omitted from any retained span, so the package carries no omission record.  No retained line contains a citation, so the wrapper carries no bibliography and no bibliography entry.  No retained line contains a figure environment, an image-inclusion command or drawing code, so no image is delivered and the package declares no asset dependency.  Editorial formatting only, in addition: an AMS \texttt{amsart} preamble that restates the article's own declarations and macros used by the retained lines, namely the environments \texttt{lem}, \texttt{notation}, \texttt{prop} on separate counters, together with the standard packages those lines need.  The retained environments print in this document as Lemma 1, Notation 1, Proposition 1 in order of appearance, whereas the article prints the same items with the numbering of its own full document.  The complete native source of the article as distributed in the arXiv e-print for this exact version is bundled unchanged as \texttt{sources/195815/source0.tex}.
\begin{lem}\label{lem:ind-join}
For finite graphs $G_1$ and $G_2$ there is a natural isomorphism
\[
  \Delta(G_1\sqcup G_2)\ \cong\ \Delta(G_1)*\Delta(G_2),
\]
where $*$ denotes the simplicial join.
In particular, $\Delta(K_2)\cong S^0$, and if $G$ is a disjoint union of $k$ edges then
\[
  \Delta(G)\ \cong\ (S^0)^{*k}\ \simeq\ S^{k-1}.
\]
\end{lem}

\begin{lem}
\label{lem:restriction-matching}
Let $Y\subseteq X$ be a subcomplex and let $M$ be an acyclic matching on $X$. Consider the matching
$M|_Y$ obtained by restricting $M$ to those matched pairs whose two faces both lie in $Y$. Then
$M|_Y$ is an acyclic matching on~$Y$.

In particular, every face of $Y$ that is critical for $M$ is also critical for $M|_Y$.
\end{lem}

\begin{notation}
\label{setup:cycle}
Let $H$ denote the induced subgraph of $C_{3k}$ on $W := V(C_{3k})\setminus \mathcal{C}_1$. The edges of $H$ are $\{x_{3i-1},x_{3i}\}$ for $i=1,\ldots k$. So, $H$ is a disjoint union of $k$ edges.


Consider the independence complexes
\[
  \Delta := \Delta(C_{3k}), \qquad A := \Delta\big(H\big),
\]
and the inclusion of simplicial complexes $i\colon A\hookrightarrow \Delta$.

Let
\[
  \sigma_1 :=  \{2,5,8,\dots,3k-1\},\qquad
  \sigma_2 :=  \{3,6,\dots,3k\}.
\]
Notice  $\sigma_1$ and $\sigma_2$ are facets of $A\subset \Delta$
 since   $\mathcal{C}_2$ and $\mathcal{C}_3$ are maximal independent sets in $H$ and $C_{3k}$.  
 \end{notation}

\begin{lem}
\label{lem:Delta-critcells}
In the situation of Setup~\ref{setup:cycle}, we have
\[\Delta  \simeq S^{k-1} \vee S^{k-1}.\]
Moreover, there exists an acyclic matching $M_\Delta$ on the face poset of $\Delta$ such that:
\begin{enumerate}
  \item $M_\Delta$ has no critical simplices of dimension $\ge k$; and
  \item the only critical $(k-1)$-simplices are exactly $\sigma_1$ and $\sigma_2$.
\end{enumerate}
Consequently $\widetilde H_{k-1}(\Delta;\Bbbk)\cong \Bbbk^2$, with basis represented by the classes of $\sigma_1$ and $\sigma_2$.
\end{lem}

\begin{prop}
\label{prop:inclusion-nonzero}
In the situation of Setup~\ref{setup:cycle}, the inclusion
\[
  i\colon A\hookrightarrow\Delta
\]
induces a nonzero homomorphism
\[
  i_{\star}: \widetilde H_{k-1}(A;\Bbbk) \longrightarrow \widetilde H_{k-1}(\Delta;\Bbbk).
\]
In particular, since $\dim_\Bbbk \widetilde H_{k-1}(A;\Bbbk)=1$, the map $i_{\star}$ is injective.
\end{prop}

\begin{proof}
Restrict the matching $M_\Delta$ from Lemma~\ref{lem:Delta-critcells} to the subcomplex $A$ to obtain a matching $M_A$ on $A$. By Lemma~\ref{lem:restriction-matching}, $M_A$ is an acyclic matching on $A$. The faces $\sigma_1$ and $\sigma_2$ lie in $A$ and are unmatched in $M_\Delta$. Hence, they are also critical $(k-1)$-simplices for $M_A$. In addition, note that $M_A$ has no critical simplices of dimension $\geq k$ since $M_{\Delta}$ does not have any such critical cells.

Let $C^{M_A}_\ast(A)$ and $C^{M_\Delta}_\ast(\Delta)$ be the Morse complexes associated to $M_A$ and $M_\Delta$, respectively. In degree $k-1$, we have
\[
  C^{M_A}_{k-1}(A) \;\cong\; \Bbbk\cdot\sigma_1 \;\oplus\; \Bbbk\cdot\sigma_2,
\]
because these are exactly the critical $(k-1)$--faces for $M_A$. The Morse differential
\[
  \partial_{k-1}^{M_A}\colon C^{M_A}_{k-1}(A) \longrightarrow C^{M_A}_{k-2}(A)
\]
is some $\Bbbk$--linear map determined by gradient paths. Since $\widetilde H_{k-1}(A;\Bbbk)\cong \Bbbk$ by \Cref{lem:ind-join} and $C^{M_A}_{k}(A)=0$, we have
\[
  \dim_\Bbbk \ker \partial_{k-1}^{M_A} = 1.
\]
So, there exist $a,b\in\Bbbk$ with $(a,b)\neq (0,0)$ such that
\[
  \partial_{k-1}^{M_A}(a\sigma_1 + b\sigma_2) = 0.
\]
Thus, $\ker \partial_{k-1}^{M_A}$ is spanned by $a\sigma_1 + b\sigma_2$. The corresponding homology class
\[
  [\omega] := [ a\sigma_1 + b\sigma_2 ] \in \widetilde H_{k-1}(A;\Bbbk)
\]
is a generator of $\widetilde H_{k-1}(A;\Bbbk)$.


Since $M_\Delta$ has no critical simplices in dimension $\ge k$, we have
$C^{M_\Delta}_{k}( \Delta)=0$. Moreover,
$C^{M_\Delta}_{k-1}(\Delta)\cong \Bbbk\cdot\sigma_1\oplus \Bbbk\cdot\sigma_2$
and $\widetilde H_{k-1}(\Delta;\Bbbk)\cong \Bbbk^2$ by Lemma~\ref{lem:Delta-critcells}.
Therefore $\partial^{M_\Delta}_{k-1}=0$, and the classes of $\sigma_1$ and $\sigma_2$
form a basis of $\widetilde H_{k-1}(\Delta;\Bbbk)$.



The inclusion $i\colon A\hookrightarrow\Delta$ induces a chain map between Morse complexes
\[
  i_{\star}\colon C^{M_A}_\ast(A) \longrightarrow C^{M_\Delta}_\ast(\Delta),
\]
which in degree $k-1$ is  the inclusion
\[
  i_{\star}\colon \Bbbk\cdot\sigma_1 \oplus \Bbbk\cdot\sigma_2
  \hookrightarrow 
  \Bbbk\cdot\sigma_1 \oplus \Bbbk\cdot\sigma_2,
\]
sending $\sigma_j$ to the same simplex viewed in $\Delta$, for $j=1,2$. So,  we obtain
\[
  i_{\star}([\omega]) 
  \;=\; [ a\sigma_1 + b\sigma_2 ] 
  \;\in\; \widetilde H_{k-1}(\Delta;\Bbbk)
  \;\cong\; \Bbbk^2.
\]
Since $[\sigma_1],[\sigma_2]$ is a basis of $\widetilde H_{k-1}(\Delta;\Bbbk)$ and $(a,b)\neq (0,0)$, we have $i_{\star}([\omega])\neq 0$. Therefore, $i_{\star}$ is a nonzero  map on $\widetilde H_{k-1}$ and it is injective.
\end{proof}

\end{document}
